% sec3_1.tex — Section 3.1: Endogeneity Bias: Working's Example

%% ─── A Simultaneous Equations Model of Market Equilibrium ──────────────────
\subsection*{A Simultaneous Equations Model of Market Equilibrium}

The classic illustration of biases created by endogenous regressors is Working
(1927).  Consider the following simple model of demand and supply:
\begin{align}
  q_i^d &= \alpha_0 + \alpha_1 p_i + u_i, \quad \text{(demand equation)}
  \label{eq:3.1.1a} \\
  q_i^s &= \beta_0 + \beta_1 p_i + v_i, \quad \text{(supply equation)}
  \label{eq:3.1.1b} \\
  q_i^d &= q_i^s, \quad \text{(market equilibrium)}
  \label{eq:3.1.1c}
\end{align}
where $q_i^d$ is the quantity demanded for the commodity in question (coffee, say) in
period $i$, $q_i^s$ is the quantity supplied, and $p_i$ is the price.  The error term $u_i$
in the demand equation represents factors that influence coffee demand other than price,
such as the public's mood for coffee.  Depending on the value of $u_i$, the demand
curve in the price-quantity plane shifts up or down.  Similarly, $v_i$ represents supply
factors other than price.  We assume that $\E(u_i) = 0$ and $\E(v_i) = 0$ (if not, include
the nonzero means in the intercepts $\alpha_0$ and $\beta_0$).  To avoid inessential complications,
we also assume $\Cov(u_i, v_i) = 0$.  If we define $q_i = q_i^d = q_i^s$, the three-equation
system \eqref{eq:3.1.1a}--\eqref{eq:3.1.1c} can be reduced to a two-equation system:
\begin{align}
  q_i &= \alpha_0 + \alpha_1 p_i + u_i, \quad \text{(demand equation)}
  \label{eq:3.1.2a} \\
  q_i &= \beta_0 + \beta_1 p_i + v_i. \quad \text{(supply equation)}
  \label{eq:3.1.2b}
\end{align}

We say that a regressor is \textbf{endogenous} if it is not predetermined (i.e., not
orthogonal to the error term), that is, if it does not satisfy the orthogonality condition.
When the equation includes the intercept, the orthogonality condition is violated and
hence the regressor is endogenous, if and only if the regressor is correlated with the
error term.  In the present example, the regressor $p_i$ is necessarily endogenous in both
equations.  To see why, treat \eqref{eq:3.1.2a}--\eqref{eq:3.1.2b} as a system of two
simultaneous equations and solve for $(p_i, q_i)$ as
\begin{align}
  p_i &= \frac{\beta_0 - \alpha_0}{\alpha_1 - \beta_1}
        + \frac{v_i - u_i}{\alpha_1 - \beta_1},
  \label{eq:3.1.3a} \\
  q_i &= \frac{\alpha_1 \beta_0 - \alpha_0 \beta_1}{\alpha_1 - \beta_1}
        + \frac{\alpha_1 v_i - \beta_1 u_i}{\alpha_1 - \beta_1}.
  \label{eq:3.1.3b}
\end{align}
So price is a function of the two error terms.  From \eqref{eq:3.1.3a}, we can calculate the
covariance of the regressor $p_i$ with the demand shifter $u_i$ and the supply shifter $v_i$
to be
\begin{equation}
  \Cov(p_i, u_i) = -\frac{\Var(u_i)}{\alpha_1 - \beta_1}, \quad
  \Cov(p_i, v_i) = \frac{\Var(v_i)}{\alpha_1 - \beta_1},
  \label{eq:3.1.4}
\end{equation}
which are not zero (unless $\Var(u_i) = 0$ and $\Var(v_i) = 0$).  Therefore, price is
correlated positively with the demand shifter and negatively with the supply shifter,
if the demand curve is downward-sloping ($\alpha_1 < 0$) and the supply curve upward-sloping
($\beta_1 > 0$).  In this example, endogeneity is a result of market equilibrium.

%% ─── Endogeneity Bias ─────────────────────────────────────────────────────
\subsection*{Endogeneity Bias}

When quantity is regressed on a constant and price, does it estimate the demand
curve or the supply curve?  The answer is neither, because price is endogenous in
both the demand and supply equations.  Recall from Section~2.9 that the OLS estimator
is consistent for the least squares projection coefficients.  In the least squares
projection of $q_i$ on a constant and $p_i$, the coefficient of $p_i$ is
$\Cov(p_i, q_i)/\Var(p_i)$,\footnote{Fact: Let $\gamma$ be the least squares coefficients
in $\tilde{\E}^*(y \mid 1, \mathbf{x}) = \alpha + \mathbf{x}'\gamma$.  Then
$\gamma = \Var(\mathbf{x})^{-1}\Cov(\mathbf{x}, y)$.  Proving this was an analytical
exercise for Chapter~2.} so,
\begin{equation}
  \plim \text{ of the OLS estimate of the price coefficient}
  = \frac{\Cov(p_i, q_i)}{\Var(p_i)}.
  \label{eq:3.1.5}
\end{equation}
To rewrite this ratio in relation to the price effect in the demand curve ($\alpha_1$), use the
demand equation \eqref{eq:3.1.2a} to calculate $\Cov(p_i, q_i)$ as
\begin{equation}
  \Cov(p_i, q_i) = \alpha_1 \Var(p_i) + \Cov(p_i, u_i).
  \label{eq:3.1.6}
\end{equation}
Substituting \eqref{eq:3.1.6} into \eqref{eq:3.1.5}, we obtain the expression for the
asymptotic bias for $\alpha_1$
\begin{equation}
  \plim \text{ of the OLS estimate of the price coefficient}
  - \alpha_1
  = \frac{\Cov(p_i, u_i)}{\Var(p_i)}.
  \label{eq:3.1.7}
\end{equation}
Similarly, the asymptotic bias for $\beta_1$, the price effect in the supply curve, is
\begin{equation}
  \plim \text{ of the OLS estimate of the price coefficient}
  - \beta_1
  = \frac{\Cov(p_i, v_i)}{\Var(p_i)}.
  \label{eq:3.1.8}
\end{equation}
But, as seen from \eqref{eq:3.1.4}, $\Cov(p_i, u_i) \neq 0$ and $\Cov(p_i, v_i) \neq 0$,
so the OLS estimate is consistent for neither $\alpha_1$ nor $\beta_1$.  This phenomenon is known as
the \textbf{endogeneity bias}.  It is also known as the \textbf{simultaneous equations bias}
or \textbf{simultaneity bias}, because the regressor and the error term are often related to
each other through a system of simultaneous equations, as in the present example.

In the extreme case of no demand shifters (so that $u_i = 0$ for all $i$), we have
$\Cov(p_i, u_i) = 0$, and the formula \eqref{eq:3.1.7} indicates that the OLS estimate is
consistent for the demand parameter $\alpha_1$.  In this case, the demand curve does not shift,
and, as illustrated in Figure~3.1(a), all the observed combinations of price and quantity
fall on the demand curve as the supply curve shifts.  In the other extreme case of no supply
shifters, the observed price-quantity pairs trace out the supply curve (see
Figure~3.1(b)).  In the general case of both curves shifting, the OLS estimate should be
consistent for a weighted average of $\alpha_1$ and $\beta_1$.  This can be seen analytically by
deriving yet another expression for the plim of the OLS estimate:
\begin{equation}
  \plim \text{ of the OLS estimate of the price coefficient}
  = \frac{\alpha_1 \, \Var(v_i) + \beta_1 \, \Var(u_i)}{\Var(v_i) + \Var(u_i)}.
  \label{eq:3.1.9}
\end{equation}
Proving this is a review question.

%% ─── Figure 3.1 placeholder ─────────────────────────────────────────────
\begin{figure}[htbp]
  \centering
  % TODO: Replace with actual figure
  \includegraphics[width=0.9\textwidth,trim=50 350 50 350,clip]{figures/fig_3_1.png}
  \caption{Supply and Demand Curves}
  \label{fig:3.1}
\end{figure}

%% ─── Observable Supply Shifters ────────────────────────────────────────────
\subsection*{Observable Supply Shifters}

The reason neither the demand curve nor the supply curve is consistently estimated
in the general case is that we cannot infer from data whether the change in price
and quantity is due to a demand shift or a supply shift.  This suggests that it might
be possible to estimate the demand curve (resp.\ the supply curve) if some of the
factors that shift the supply curve (resp.\ the demand curve) are observable.  So
suppose the supply shifter, $v_i$, can be divided into an observable factor $x_i$ and an
unobservable factor $\zeta_i$ uncorrelated with $x_i$.\footnote{This decomposition is
always possible.  If the least squares projection of $v_i$ on a constant and $x_i$ is
$\gamma_0 + \beta_2 x_i$, define $\zeta_i \equiv v_i - \gamma_0 - \beta_2 x_i$, so
$v_i = \zeta_i + \gamma_0 + \beta_2 x_i$.  By definition, $\zeta_i$ is uncorrelated with $x_i$.
Substitute this equation into the original supply equation \eqref{eq:3.1.2b} and submerge
$\gamma_0$ in the intercept.  This produces \eqref{eq:3.1.2bp}.}
\begin{equation}
  q_i = \beta_0 + \beta_1 p_i + \beta_2 x_i + \zeta_i
  \quad \text{with } \beta_2 \neq 0. \quad \text{(supply)}
  \label{eq:3.1.2bp}
\end{equation}
Now imagine that this observed supply shifter $x_i$ is predetermined (i.e., uncorrelated
with the error term) in the demand equation; think of $x_i$ as the temperature
in coffee-growing regions.  If the temperature ($x_i$) is uncorrelated with the unobserved
taste for coffee ($u_i$), it should be possible to extract from price movements
a component that is related to the temperature (the observed supply shifter) but
uncorrelated with the demand shifter.  We can then estimate the demand curve
by examining the relationship between coffee consumption and that component of
price.  Let us formalize this argument.

For the equation in question, a predetermined variable that is correlated with
the endogenous regressor is called an \textbf{instrumental variable} or an \textbf{instrument}.
We sometimes call it a \textbf{valid instrument} to emphasize that its correlation with the
endogenous regressor is not zero.  In this example, the observable supply shifter $x_i$
can serve as an instrument for the demand equation.  This can be seen easily.  Solve
the system of two simultaneous equations \eqref{eq:3.1.2a} and \eqref{eq:3.1.2bp} for
$(p_i, q_i)$:
\begin{align}
  p_i &= \frac{\beta_0 - \alpha_0}{\alpha_1 - \beta_1}
        + \frac{\beta_2}{\alpha_1 - \beta_1}\, x_i
        + \frac{\zeta_i - u_i}{\alpha_1 - \beta_1},
  \label{eq:3.1.10a} \\
  q_i &= \frac{\alpha_1 \beta_0 - \alpha_0 \beta_1}{\alpha_1 - \beta_1}
        + \frac{\alpha_1 \beta_2}{\alpha_1 - \beta_1}\, x_i
        + \frac{\alpha_1 \zeta_i - \beta_1 u_i}{\alpha_1 - \beta_1}.
  \label{eq:3.1.10b}
\end{align}
Since $\Cov(x_i, \zeta_i) = 0$ by construction and $\Cov(x_i, u_i) = 0$ by assumption, it
follows from \eqref{eq:3.1.10a} that
\[
  \Cov(x_i, p_i) = \frac{\beta_2}{\alpha_1 - \beta_1}\, \Var(x_i) \neq 0.
\]
So $x_i$ is indeed a valid instrument.

With a valid instrument at hand, we can estimate the price coefficient $\alpha_1$ of
the demand curve consistently.  Use the demand equation \eqref{eq:3.1.2a} to calculate, not
$\Cov(p_i, q_i)$ as in \eqref{eq:3.1.6}, but $\Cov(x_i, q_i)$:
\begin{align*}
  \Cov(x_i, q_i) &= \alpha_1 \Cov(x_i, p_i) + \Cov(x_i, u_i) \\
  &= \alpha_1 \Cov(x_i, p_i) \quad (\Cov(x_i, u_i) = 0 \text{ by assumption}).
\end{align*}
As we just verified, $\Cov(x_i, p_i) \neq 0$.  So we can divide both sides by $\Cov(x_i, p_i)$
to obtain
\begin{equation}
  \alpha_1 = \frac{\Cov(x_i, q_i)}{\Cov(x_i, p_i)}.
  \label{eq:3.1.11}
\end{equation}
A natural estimator that suggests itself is
\begin{equation}
  \hat{\alpha}_{1,\text{IV}}
  = \frac{\text{sample covariance between } x_i \text{ and } q_i}
         {\text{sample covariance between } x_i \text{ and } p_i}.
  \label{eq:3.1.12}
\end{equation}
This estimator is called the \textbf{instrumental variables (IV) estimator} with $x_i$ as the
instrument.  We sometimes say ``the endogenous regressor $p_i$ is instrumented by $x_i$.''

Another popular procedure for consistently estimating $\alpha_1$ is the \textbf{two-stage least
squares (2SLS)}.  It is so called because the procedure consists of running two
regressions.  In the first stage, the endogenous regressor $p_i$ is regressed on a constant
and the predetermined variable $x_i$, to obtain the fitted value, $\hat{p}_i$, of $p_i$.  The
second stage is to regress the dependent variable $q_i$ on a constant and $\hat{p}_i$.  The use of
$\hat{p}_i$ rather than $p_i$ distinguishes 2SLS from the simple-minded application of OLS
to the demand equation.  The 2SLS estimator of $\alpha_1$ is the OLS estimate of the $\hat{p}_i$
coefficient in the second-stage regression.  So it equals\footnote{Fact: In the regression
of $y_i$ on a constant and $x_i$, the OLS estimator of the $x_i$ coefficient is the ratio of the
sample covariance between $x_i$ and $y_i$ to the sample variance of $x_i$.  Proving this was a
review question of Section~1.2.}
\begin{equation}
  \hat{\alpha}_{1,\text{2SLS}}
  = \frac{\text{sample covariance between } \hat{p}_i \text{ and } q_i}
         {\text{sample variance of } \hat{p}_i}.
  \label{eq:3.1.13}
\end{equation}

To relate the regression in the second stage to the demand equation, rewrite
\eqref{eq:3.1.2a} as
\begin{equation}
  q_i = \alpha_0 + \alpha_1 \hat{p}_i + [u_i + \alpha_1 (p_i - \hat{p}_i)].
  \label{eq:3.1.14}
\end{equation}
The second stage regression estimates this equation, treating the bracketed term as
the error term.  The OLS estimate of $\alpha_1$ is consistent for the following reason.  If the
fitted value $\hat{p}_i$ were exactly equal to the least squares projection
$\tilde{\E}^*(p_i \mid 1, x_i)$, then
neither $u_i$ nor $(p_i - \hat{p}_i)$ would be correlated with $\hat{p}_i$: $u_i$ is uncorrelated because it is
uncorrelated with $x_i$ and $\hat{p}_i$ is a linear function of $x_i$, and $(p_i - \hat{p}_i)$ is uncorrelated
because it is a least squares projection error.  The fitted value $\hat{p}_i$ is not exactly equal
to $\tilde{\E}^*(p_i \mid 1, x_i)$, but the difference between the two vanishes as the sample gets
larger.  Therefore, asymptotically, $\hat{p}_i$ is uncorrelated with the error term in
\eqref{eq:3.1.14}, making the 2SLS estimator consistent.  Furthermore, since the projection
is the best linear predictor, the fitted value incorporates all of the effect of the supply
shifter on price.  This suggests that minimizing the asymptotic variance is minimized for
the 2SLS estimator.

In the present example, the IV and 2SLS estimators are numerically the same
(we will prove this in a more general context).  More generally, the 2SLS estimator
can be written as an IV estimator with an appropriate choice of instruments, and
the IV estimator, in turn, is a particular GMM estimator.

%% ─── Questions for Review ──────────────────────────────────────────────────
\bigskip
\begin{center}\rule{0.6\textwidth}{0.4pt}\end{center}
\noindent\textsc{Questions for Review}
\medskip

\begin{enumerate}
\item In the simultaneous equations model \eqref{eq:3.1.2a}--\eqref{eq:3.1.2b},
  suppose $\Cov(u_i, v_i)$ is not necessarily zero.  Is it necessarily true that price and the
  demand shifter $u_i$ are positively correlated when $\alpha_1 < 0$ and $\beta_1 > 0$?
  [Answer: No.]  Why?
\end{enumerate}
